\documentclass[10pt,a4paper]{amsart}
\usepackage[margin=19mm]{geometry}
\usepackage{amsmath,amssymb,mathtools}
\usepackage[scr=rsfs,cal=euler]{mathalfa}
\usepackage{xfrac}
\usepackage[unicode,hidelinks,bookmarks=false]{hyperref}
\newcommand{\ie}{i.e.\ }
\newcommand{\cf}{cf.\ }
\newcommand{\resp}{respectively\ }
\newcommand{\Subsection}{\subsection}
\providecommand{\nobreakdash}{\nobreak\hbox{-}}
\setlength{\emergencystretch}{2em}
\multlinegap0pt
\usepackage{amssymb}
\usepackage{bm}
\newtheorem{theorem}{Theorem}
\title{A Family of Generating Functions for Reciprocal Binomial Coefficients and Its Applications}
\author{Dmitry Kruchinin, Vladimir Kruchinin}
\date{Source: arXiv:2602.07822v1 (CC BY 4.0)}
\begin{document}
\maketitle

\noindent\emph{Source and licence.} A Family of Generating Functions for Reciprocal Binomial Coefficients and Its Applications. Version arXiv:2602.07822v1 (CC BY 4.0), distributed by arXiv under the Creative Commons Attribution 4.0 International licence, http://creativecommons.org/licenses/by/4.0/, as recorded in the arXiv licence metadata for this exact version and shown on the abstract page https://arxiv.org/abs/2602.07822v1. Full licence notice: \texttt{licenses/arxiv-2602.07822-CC-BY-4.0.txt}.\par\smallskip

\noindent\textit{Editorial scope.} Counted unit: the article's theorem iden6 (source0.tex lines 911--918). The article's complete original proof of it is delivered with it (source0.tex lines 920--943, one \texttt{proof} environment of 24 lines). No mathematics is written, repaired or re-derived: the statement and the proof are the article's own lines, in the article's own notation, and no retained line is substituted or re-flowed.  No further source-local context is needed: every cross-reference of every retained line resolves inside the two delivered spans.  Nothing was omitted from any retained span, so the package carries no omission record.  No retained line contains a citation, so the wrapper carries no bibliography and no bibliography entry.  No retained line contains a figure environment, an image-inclusion command or drawing code, so no image is delivered and the package declares no asset dependency.  Editorial formatting only, in addition: an AMS \texttt{amsart} preamble that restates the article's own declarations and macros used by the retained lines, namely the environments \texttt{theorem} on separate counters, together with the standard packages those lines need.  The retained environments print in this document as Theorem 1 in order of appearance, whereas the article prints the same items with the numbering of its own full document.  The complete native source of the article as distributed in the arXiv e-print for this exact version is bundled unchanged as \texttt{sources/194215/source0.tex}.
\begin{theorem}
 For \( n \geq 1 \) the following holds:
\begin{equation}\label{iden6}
\sum_{m=0}^{\lfloor (n-1)/2 \rfloor} \frac{1}{(n-m)\binom{n-m-1}{m}}
= \sum_{i=1}^{n} \frac{F_i\bigl((-1)^n + 2(-1)^{i-1}\bigr)}{n-i+1},
\end{equation}
where \( F_i \) is the \( i \)-th Fibonacci number (\( F_1 = F_2 = 1 \)).
\end{theorem}

\begin{proof}
 Consider the diagonal sum for \( I(x,y) \):
\[
S_{10}(x) = I(x,x) = \sum_{n\geq 1}\left(
\sum_{m=0}^{\lfloor (n-1)/2 \rfloor} \frac{1}{(n-m)\binom{n-m-1}{m}}
\right)x^n.
\]
Using the explicit form of \( I(x,x) \), we get:
\[
I(x,x) = \frac{\log\!\bigl((1-x)(1-x^2)\bigr)}{x^2+x-1}
= \frac{2\log(1-x)}{1+x-x^2} - \frac{\log(1+x)}{1+x-x^2}.
\]
Note that
\[
\frac{1}{1+x-x^2} = \sum_{n\geq 0} (-1)^n F_{n+1} x^n.
\]
Multiplication by \( \log(1\pm x) \) corresponds to taking a convolution:
\[
\frac{\log(1-x)}{1+x-x^2} = -\sum_{n\geq 1}\left(
\sum_{i=1}^n \frac{(-1)^{n-i} F_{n-i+1}}{i}
\right)x^n.
\]
Combining terms, we arrive at formula \eqref{iden6}.
\end{proof}

\end{document}
